\section{Coarse-Graining I: The Meta-Agent as Gauge-Covariant Barycenter}\label{sec:meta_agent_formation}

A meta-agent is the gauge-covariant coarse-graining of a coherent cluster of agents into a single agent one level up. Its anatomy, used throughout the manuscript, has four parts. The meta-agent is \emph{defined} by its state, a variational barycenter whose frame-free part is the defining datum. \emph{When} a cluster should be replaced by such a parent is the question of a free-energy licensing criterion (status C). \emph{How} candidate clusters are found is a threshold detector, a coherence filter (status S) not derived from the criterion. Finally, gauge covariance makes the barycenter a frame-consistent coarse-graining map, not an exact one. The companion theory supersedes both the barycenter and the detector: its parents are exact pushforwards of the fine laws under declared coarse channels~\cite[Chs.~9 and~11]{DennisGaugeVFE2026}, and its formation results are witnesses for declared families~\cite[Ch.~13]{DennisGaugeVFE2026}. The barycenter is retained here as a working construction; its relation to the exact parent is stated at the end of Section~\ref{sec:meta_agent_definition}, and the exact coarse-graining step of this manuscript is developed in Section~\ref{sec:rg_status} and Appendix~\ref{app:rigorous_rg}.

\subsection{The Meta-Agent as Gauge-Covariant Variational Barycenter}
\label{sec:meta_agent_definition}

The meta-agent's belief is the M-projection barycenter of its transported constituents. Conditional on a parent existing, treating the parent as a variational distribution and minimizing the alignment energy along the child-parent edges over $q_I$ gives
\begin{equation}
q_I^* = \arg\min_{q_I} \sum_{i \in I} w_i^I D_{\mathrm{KL}}(\Omega_{Ii} q_i \| q_I) + \lambda D_{\mathrm{KL}}(q_I \| p_I),
\label{eq:meta_agent_barycenter}
\end{equation}
a moment-matching barycenter, with the meta-belief $q_I$ in the second slot, regularized toward a prior $p_I$ on the parent. The belief coupling $D_{\mathrm{KL}}(q_i \| \Omega_{ij} q_j)$ of Eq.~\eqref{eq:beta_optimal} places the receiver's law in the first slot (the I-projection direction), so the two constructions use opposite slot orders. Here $w_i^I \ge 0$ is the cluster-aggregation weight of constituent $i$ in cluster $I$; it is a distinct object from the per-agent precision $\alpha_i$ of Section~\ref{sec:state_dependent_precision}, which weights the agent's own belief-prior coupling $D_{\mathrm{KL}}(q_i \| p_i)$. For Gaussian beliefs $q_i = \mathcal{N}(\mu_i, \Sigma_i)$ with transported moments $\tilde{\mu}_i = \Omega_{Ii}\mu_i$ and $\tilde{\Sigma}_i = \Omega_{Ii}\Sigma_i\Omega_{Ii}^\top$, the closed form (with $W_I := \sum_{i \in I} w_i^I$ and $\lambda \to 0$ for legibility) is
\begin{align}
\mu_I^* &= \tfrac{1}{W_I}\textstyle\sum_i w_i^I \tilde{\mu}_i,
\label{eq:meta_agent_mu_barycenter}
\\
\Sigma_I^* &= \tfrac{1}{W_I}\textstyle\sum_i w_i^I\big[\tilde{\Sigma}_i + (\mu_I^* - \tilde{\mu}_i)(\mu_I^* - \tilde{\mu}_i)^\top\big].
\label{eq:meta_agent_sigma_barycenter}
\end{align}
The mean is the $w^I$-weighted average of the transported constituent means; the covariance is the $w^I$-weighted average of the transported covariances plus a dispersion term that records how much the constituents disagree about the mean. The dispersion term vanishes at strict consensus, when all $\tilde{\mu}_i$ coincide. In the high-coherence regime it is $\mathcal{O}(\varepsilon)$ in the post-transport pairwise dispersion, and dropping it is an approximation, not an identity (the error bound and the working forms are given in Section~\ref{sec:meta_agent_state_construction}).

The M-projection is not the only equivariant variational barycenter. The I-projection $\arg\min_q \sum_i w_i^I D_{\mathrm{KL}}(q \| \Omega_{Ii} q_i)$, whose Gaussian solution is the precision average $\Lambda_I = \sum_i w_i^I \tilde{\Lambda}_i / W_I$, is also a strict minimizer and is equally covariant under common pushforward. We adopt the M-projection for two reasons: it is the exact M-step coordinate of a random-effects model for the constituents (Section~\ref{sec:meta_agent_variational}), and it pools the constituents' uncertainty, whereas the precision average sharpens it.

\begin{proposition}[Existence, uniqueness, and equivariance of the belief barycenter]
\label{prop:barycenter_existence}
Let $w_i^I \ge 0$ with $W_I = \sum_{i \in I} w_i^I > 0$, and let the transported covariances $\tilde{\Sigma}_i = \Omega_{Ii}\Sigma_i\Omega_{Ii}^\top$ be positive definite. Then the minimizer of Eq.~\eqref{eq:meta_agent_barycenter} with $\lambda \to 0$ over the Gaussian family exists and is unique, given by Eqs.~\eqref{eq:meta_agent_mu_barycenter}--\eqref{eq:meta_agent_sigma_barycenter}, and the map $(\{q_i\}, \{U_i\}, U_I) \mapsto q_I^*$ commutes with the action of one global $g \in G$ on beliefs and frames (Theorem~\ref{thm:rg_covariance}).
\end{proposition}

\begin{proof}
In the natural parameters $\theta$ of $q_I$, each term $D_{\mathrm{KL}}(\Omega_{Ii}q_i \| q_\theta)$ is, up to a $\theta$-independent constant, the Bregman divergence $B_\psi(\theta, \cdot)$ of the Gaussian log-partition function $\psi$, which is strictly convex on the open natural-parameter domain, so the objective equals $\sum_{i \in I} w_i^I[\psi(\theta) - \langle \tilde{\eta}_i, \theta\rangle] + \mathrm{const}$ with $\tilde{\eta}_i$ the expectation parameter of the transported constituent $\Omega_{Ii}q_i$. Strict convexity gives a unique stationary point $\nabla\psi(\theta^*) = W_I^{-1}\sum_{i \in I} w_i^I \tilde{\eta}_i$, which is exactly the moment matching of Eqs.~\eqref{eq:meta_agent_mu_barycenter}--\eqref{eq:meta_agent_sigma_barycenter}, and the minimizer exists because a convex combination of interior expectation parameters lies in the open expectation domain. Equivariance is the statement of Theorem~\ref{thm:rg_covariance}.
\end{proof}

The same closed form, together with its minimized score $\tfrac12[\log\det\Sigma_I^* - \sum_i w_i^I \log\det\tilde{\Sigma}_i]$ for normalized weights, is the full-Gaussian barycenter proposition of~\cite[Ch.~G8]{DennisGaugeVFESuppG2026}. For $\lambda > 0$ the regularized objective adds the term $\lambda D_{\mathrm{KL}}(q_I \| p_I)$, with the parent in the first slot, for which coercivity yields existence; uniqueness is claimed only for the $\lambda \to 0$ statement above.

The parent's gauge frame requires separate treatment, because it is not identifiable from anything the parent does. The next proposition states this and a companion obstruction.

\begin{proposition}[The parent frame is pure gauge; no symmetric equivariant frame mean]
\label{prop:parent_frame_gauge}
(i) Let a parent with frame $U_I$ and law $\mathcal{N}(\mu_I, \Sigma_I)$ in that frame supply constituent $i$ the transported law $(\Omega_{iI})_\#\mathcal{N}(\mu_I, \Sigma_I)$ with $\Omega_{iI} = U_i U_I^{-1}$, as in the cross-scale prior of Eq.~\eqref{eq:cross_scale_shadow}. This law is $\mathcal{N}(U_i m, U_i S U_i^\top)$ with
\begin{equation*}
m = U_I^{-1}\mu_I, \qquad S = U_I^{-1}\Sigma_I U_I^{-\top},
\end{equation*}
so every constituent law depends on $(U_I, \mu_I, \Sigma_I)$ only through the frame-free pair $(m, S)$. The barycenter of Eqs.~\eqref{eq:meta_agent_mu_barycenter}--\eqref{eq:meta_agent_sigma_barycenter}, computed with $\Omega_{Ii} = U_I U_i^{-1}$ and weights built from transported divergences, gives a pair $(m, S)$ that does not depend on $U_I$.
(ii) For $K_q \ge 2$ and $n \ge 2$ there is no map $M: \mathrm{GL}^+(K_q)^n \to \mathrm{GL}^+(K_q)$ that is symmetric under permutations of its arguments and left equivariant, $M(gU_1, \ldots, gU_n) = g M(U_1, \ldots, U_n)$ for all $g$; no continuity is assumed, and the same holds with $\mathrm{SO}(K_q)$ in place of $\mathrm{GL}^+(K_q)$.
\end{proposition}

\begin{proof}
(i) The pushforward of $\mathcal{N}(\mu_I, \Sigma_I)$ by $\Omega_{iI}$ is $\mathcal{N}(U_i U_I^{-1}\mu_I, U_i U_I^{-1}\Sigma_I U_I^{-\top} U_i^\top)$, which is the stated law. For the barycenter, $\tilde{\mu}_i = U_I (U_i^{-1}\mu_i)$ and $\tilde{\Sigma}_i = U_I (U_i^{-1}\Sigma_i U_i^{-\top}) U_I^\top$, so Eqs.~\eqref{eq:meta_agent_mu_barycenter}--\eqref{eq:meta_agent_sigma_barycenter} give $m = W_I^{-1}\sum_i w_i^I U_i^{-1}\mu_i$ and $S = W_I^{-1}\sum_i w_i^I [U_i^{-1}\Sigma_i U_i^{-\top} + (m - U_i^{-1}\mu_i)(m - U_i^{-1}\mu_i)^\top]$, in which $U_I$ does not appear; transported divergences are invariant under the common pushforward by $U_I^{-1}$, so the weights do not depend on $U_I$ either. (ii) Let $v \in \mathrm{SO}(K_q)$ be a rotation by $2\pi/n$ in one coordinate plane, so that $v$ has order $n$, and take $U_j = v^{j-1}$ for $j = 1, \ldots, n$. Left multiplication by $v$ permutes this tuple cyclically, so symmetry gives $M(vU_1, \ldots, vU_n) = M(U_1, \ldots, U_n)$, while equivariance gives $M(vU_1, \ldots, vU_n) = vM(U_1, \ldots, U_n)$. Hence $(v - I)M = 0$, and since $M$ is invertible, $v = I$, a contradiction.
\end{proof}

Part (i) says that $U_I$ ranges over a $K_q^2$-dimensional orbit along which nothing observable varies, and only $(m, S)$ is determined. A search for the correct parent frame is therefore a search for a canonical value of a pure-gauge quantity: any rule fixing $U_I$ is a gauge convention to be declared, and the open problem recorded in Section~\ref{sec:open_problems} concerns the choice of such a convention, not a theorem. Part (ii) is a global statement. Frame means that are symmetric and equivariant exist on restricted domains of sufficiently concentrated frames, which exclude the cyclic tuples of the proof.

The default convention is the Riemannian center of mass of the constituent frames,
\begin{equation}
U_I^* = \arg\min_{U \in G} \sum_{i \in I} w_i d_G(U, U_i)^2.
\label{eq:meta_agent_frame_barycenter}
\end{equation}
For a compact gauge group with bi-invariant metric $d_G(U, V)^2 = \|\log(U^{-1}V)\|_F^2$, the construction is that of Grove and Karcher~\cite{GroveKarcher1973}, with estimates by Karcher~\cite{Karcher1977} (see also~\cite{Pennec2009}). The minimizer exists, is unique as the global minimizer over $G$, and lies in the ball, whenever the weighted frames lie in an open geodesic ball of radius $\rho < r_{\mathrm{cx}} = \tfrac12\min\{\operatorname{inj} G, \pi/\sqrt{\Delta}\}$, with $\Delta$ an upper bound on the sectional curvature~\cite[Thm.~2.1]{Afsari2011}. For $\mathrm{SO}(K_q)$ with this metric, $r_{\mathrm{cx}} = \pi/\sqrt{2}$ in $d_G$, which for $\mathrm{SO}(3)$ is a rotation angle of $\pi/2$ (for $\mathrm{SO}(3)$, $\Delta = 1/8$ and $\operatorname{inj} = \sqrt{2}\pi$; for $K_q \ge 4$, $\Delta = 1/4$ with the same injectivity radius). Kendall~\cite{Kendall1990} obtains the larger radius $\min\{\operatorname{inj} G, \pi/(2\sqrt{\Delta})\}$ only for the mean defined with the distance restricted to the ball, and that radius does not govern the global argmin~\eqref{eq:meta_agent_frame_barycenter}: three rotations by $0$, $2\pi/3$, and $4\pi/3$ about one axis lie inside Kendall's radius, and the uniform-weight objective attains its minimum $16\pi^2/27$ at each of the three. On the convex ball the Karcher convention is left and right equivariant, since $d_G$ is bi-invariant. The frame condition $\rho < r_{\mathrm{cx}}$ is not detectable by the consensus score of Section~\ref{sec:meta_agent_threshold}, which is built from gauge-invariant divergences: if the three constituents above carry the beliefs $(U_j)_\# q_0$ of one law $q_0$, every transported divergence vanishes and the belief-coherence factor equals one.

The natural variational-inference gauge group $\mathrm{GL}^+(K_q)$ is noncompact and admits no bi-invariant Riemannian metric~\cite[Lemma 7.5]{Milnor1976}, a connected Lie group carrying one exactly when it is isomorphic to a product of a compact group and a vector space. The Ad-invariant trace form $\langle X, Y \rangle = \mathrm{tr}(XY)$ on $\mathfrak{gl}(K_q)$ is bi-invariant and nondegenerate, so a bi-invariant \emph{pseudo}-Riemannian metric does exist, of signature $(K_q(K_q+1)/2, K_q(K_q-1)/2)$. Being indefinite, it induces no distance function and supplies no Fr\'echet objective. The expression $\|\log(U^{-1}V)\|_F$ is not a distance on $\mathrm{GL}^+(K_q)$: the principal logarithm is undefined on an open set of the group, and the triangle inequality fails. The obstruction is to a bi-invariant variance, not to a bi-invariant mean. The exponential barycenter of the canonical Cartan-Schouten connection, the solution $U_I$ of $\sum_i w_i \log(U_I^{-1}U_i) = 0$, is left-, right- and inversion-equivariant and exists uniquely for sufficiently concentrated frames~\cite[Thms.~11--12]{PennecArsigny2013}; it is not the minimizer of any Riemannian variance. Under a left-invariant Riemannian metric instead~\cite{Milnor1976, Andruchow2014}, the group is complete and homogeneous, so Afsari's theorem applies and the Riemannian center of mass exists and is unique below $r_{\mathrm{cx}}$, but it is only left-equivariant. The affine-invariant route through the polar decomposition~\cite{PennecFillardAyache2006} is a further option. By Proposition~\ref{prop:parent_frame_gauge}, the choice among these is a choice of gauge convention with no effect on the parent laws.

Gauge covariance of the barycenter map (Theorem~\ref{thm:rg_covariance}) makes the belief barycenter and a declared frame convention frame-consistent coarse-graining maps: under the theorem's hypotheses they commute with the global action on beliefs and frames, so the parent is not an artifact of the frames in which the constituents happen to be written. Covariance does not single these maps out, since the I-projection above is equally covariant, and it does not make the barycenter an exact coarse-graining. The exact parent of a cluster is the pushforward of the fine joint law, posterior, and recognition law under one declared coarse kernel~\cite[Ch.~9]{DennisGaugeVFE2026} (the pointwise parent-datum corollary). The Gaussian barycenter with its dispersion term is the moment-matched projection of the random-selection mixture of the transported constituents~\cite[Ch.~G8]{DennisGaugeVFESuppG2026}, and the pooled form without the dispersion term, used in Section~\ref{sec:meta_agent_state_construction}, is already for two scalar constituents with equal weights the image of no parameter-independent Markov kernel~\cite[Ch.~G8]{DennisGaugeVFESuppG2026} (the normal-barycenter kernel proposition). The evidence-preservation and data-processing statements of the exact construction therefore do not transfer to the working form.

\begin{definition}[Meta-Agent]
Let $I = \{i, j, \ldots\} \subseteq \mathcal{I}^{(\ell)}$ be a cluster of level-$\ell$ agents with (1) non-empty multi-child support
\begin{equation}
\mathcal{U}_I = \mathrm{ConnComp}\Big(\big\{c \in \mathcal{C}: \textstyle\sum_{i \in I} \chi_i(c) \ge 2\big\}\Big) \neq \varnothing,
\label{eq:meta_agent_support}
\end{equation}
a connected region in which at least two constituents maintain beliefs simultaneously, with $\chi_i: \mathcal{C} \to [0, 1]$ agent $i$'s support function (introduced with the continuous-field free energy in Section~\ref{sec:complete_free_energy}; Appendix~\ref{app:rigorous_rg} takes $\chi_i \in \{0, 1\}$), and (2) minimum cluster size $|I| \geq N_{\min}$ (typically $N_{\min} = 2$). The meta-agent $\mathcal{A}^{(\ell+1)}$ of $I$ is the measurable, piecewise-smooth section $(q^{(\ell+1)}, p^{(\ell+1)}, s^{(\ell+1)}, r^{(\ell+1)}, \phi^{(\ell+1)})$ over $\mathcal{U}_I$ whose belief is the barycenter~\eqref{eq:meta_agent_barycenter} and whose model is the analogous barycenter on the model fiber. Its defining datum is the frame-free pair $(m, S)$ of Proposition~\ref{prop:parent_frame_gauge}; the frame $\phi^{(\ell+1)}$, equivalently $U_I$, is a declared gauge convention, by default the Karcher mean~\eqref{eq:meta_agent_frame_barycenter}, which is single-valued when the constituent frames at $c$ lie in a geodesic ball of radius below $r_{\mathrm{cx}}$. The averaged fields are smooth wherever the active constituent set and the weights vary smoothly, but can fail to be smooth across support boundaries where $\sum_{i \in I}\chi_i$ crosses the threshold in~\eqref{eq:meta_agent_support}, so we claim only the measurability that the coarse-graining of Appendix~\ref{app:rigorous_rg} requires. Under the smooth-bump convention of Section~\ref{sec:complete_free_energy} (smooth $\chi_i$), the averaged fields are smooth on the interior of $\mathcal{U}_I$, and it is in this interior sense that the meta-agent re-enters the free energy as an agent at level $\ell+1$. The stricter $m$-fold variant $\sum_{i \in I}\chi_i(c) \ge m$ for $m > 2$, up to the global-intersection limit $m = |I|$, is available where applications require deeper simultaneous coverage. A meta-agent is \emph{licensed} when, in addition, the criterion~\eqref{eq:meta_agent_FE_criterion} of Section~\ref{sec:meta_agent_variational} holds; the threshold detector of Section~\ref{sec:meta_agent_threshold} is a filter for candidate clusters and is not part of the definition.
\end{definition}

The slow model channel $s$ carries the persistent identity across this definition: belief consensus shifts on the perceptual timescale, while the generative model changes only on the learning timescale. On an interpretive reading (status S), a culture, a school of thought, or a scientific community is a cluster whose members share enough of a common slow generative model, after transport into one another's frames, to render their mutual disagreements intelligible. Members belong to the same meta-agent not when $s_i = s_j$, a strict identity no nontrivial culture exhibits, but when $s_i \approx \tilde\Omega_{ij} s_j$, with the transport $\tilde\Omega_{ij}$ distinguishing intra-cultural disagreement (translatable across frames) from radical incommensurability (not translatable). The fast-channel beliefs $q_i$ may diverge so long as the models remain mutually translatable. The active-inference counterpart of the construction is the group-level generative model of Waade et al.~\cite{waade2025asone}, who propose that a collective can constitute a group-level agent if it maintains a group-level Markov blanket; the detector used here is a similarity threshold and does not test that conditional-independence condition. That coherent collective structure can emerge from each agent minimizing variational free energy with no hand-coded alignment rules is the central finding of Heins et al.~\cite{heins2024collective}, the closest external precedent for the emergent reading, while the complementary fragmented limit, in which confirmation-biased free-energy-minimizing sampling produces echo chambers tunable by prior precision, is the active-inference account of Albarracin et al.~\cite{albarracin2022epistemic}. The stronger reading, on which meta-agents are dynamical agents in their own right and not statistical summaries of their constituents, is an interpretive commitment of status S (Table~\ref{tab:claim_status}); Section~\ref{sec:meta_agent_variational} shows that the exact variational reading of the barycenter treats the parent as a fitted hyperparameter.

\subsection{Licensing: The Free-Energy Improvement Criterion}
\label{sec:meta_agent_variational}

The definition above specifies what a meta-agent is, conditional on its existing. The proposed rule for when a cluster should be replaced by one is a variational free-energy improvement criterion. A candidate meta-agent $I$ is licensed when the optimized free energy with the parent inserted, plus an organizational overhead $C(I)$, falls below the optimized free energy of the disaggregated constituents,
\begin{equation}
\mathcal{F}^*[q_I, p_I, s_I, r_I, \{q_i, s_i\}_{i \in I}] + C(I) < \mathcal{F}^*[\{q_i, p_i, s_i, r_i\}_{i \in I}].
\label{eq:meta_agent_FE_criterion}
\end{equation}
Here $C(I)$ is the cost of carrying the additional organizational structure (the parent's variational parameters and its place in the cross-scale hierarchy). It is not specified here; entropy or minimum-description-length terms in the parent's primitive variables are candidates, so the criterion contains a free function. The criterion is the $\Delta_I > 0$ special case of the description-length retention rule of Appendix~\ref{sec:meta_agent_rg} (Eq.~\eqref{eq:rg_retention_gain}), where $\Delta_I$ is the net retention gain at an unspecified complexity penalty $C(I)$; the general rule, its penalty choices, and the closure ansatz under which it is admissible are deferred to Appendix~\ref{sec:meta_agent_rg}, with statements and proofs in Appendix~\ref{app:rigorous_rg}. The criterion has status C for two reasons.

First, the source of any saving is not identified. Which intra-cluster coupling terms are removed when the parent is inserted is part of a post-merge functional that the present formalism leaves unspecified. At frozen attention weights, the row constraint $\sum_j \beta_{ij} = 1$ caps the energetic intra-cluster alignment $\sum_i \sum_{j \in I} \beta_{ij} D_{\mathrm{KL}}(q_i \| \Omega_{ij} q_j)$ at $|I|\varepsilon$ when the pairwise post-transport divergences are bounded by $\varepsilon$ (not the quadratic $|I|(|I|-1)\varepsilon$ an unweighted pair count would give). This bound does not measure a saving of $\mathcal{F}^*$. At Gibbs-optimal rows each row sits at the value $A_i = -\tau\log(Z_{\mathrm{in}} + Z_{\mathrm{out}})$, with $Z_{\mathrm{in}}$ and $Z_{\mathrm{out}}$ the prior-weighted partition sums over sources inside and outside the cluster, and deleting the intra-cluster sources raises it by $\tau\log(1 + Z_{\mathrm{in}}/Z_{\mathrm{out}}) > 0$ when the prior is left unnormalized; with the prior rescaled to unit mass over the survivors the change is still positive whenever the intra-cluster energies are the smaller ones, which is the high-coherence case. Any net gain must therefore come from the rows of agents outside the cluster, which couple to one parent in place of $|I|$ constituents, or from the parent's prior term, and this is not derived here.

Second, the criterion is degenerate when the parent is exact. If the parent laws are the pushforward of the fine laws under a declared coarse channel, the fine free energy equals the coarse one plus the discarded conditional information, $\mathcal{F}_{\mathrm{fine}} = \mathcal{F}_{\mathrm{coarse}} + \Delta_A$ with $\Delta_A \ge 0$ for every channel, at unchanged evidence~\cite[Ch.~9]{DennisGaugeVFE2026} (the exact coarse free-energy theorem and its evidence-preserving corollary). Read with exact parents, Eq.~\eqref{eq:meta_agent_FE_criterion} reduces to $\Delta_A > C(I)$: it licenses whichever cluster discards the most conditional information, and the drop in free energy is neither model improvement nor evidence ascent. A usable licensing rule needs a post-merge functional that is not monotone under refinement and has a nondegenerate optimizer; neither this paper nor~\cite[Ch.~9]{DennisGaugeVFE2026}, where the selection of partitions is an open problem, supplies one.

The barycenter itself has an exact variational reading, once the generative model it belongs to is stated. Let the constituents' latent states be independent draws $k_i \sim \mathcal{N}(U_i m, U_i S U_i^\top)$ from a random-effects model whose hyperparameter $\theta = (m, S)$ is a point parameter, the frame-free parent pair of Proposition~\ref{prop:parent_frame_gauge}. In $\sum_i w_i^I \mathbb{E}_{q_i}[\log p_\theta(k_i)]$ the transport cancels from both the trace and the quadratic form, leaving a weighted Gaussian maximum-likelihood problem for the transported moments whose maximizer is exactly Eqs.~\eqref{eq:meta_agent_mu_barycenter}--\eqref{eq:meta_agent_sigma_barycenter}, dispersion term included. The barycenter is therefore the M-step, or type-II maximum-likelihood (empirical-Bayes), coordinate for $\theta$ at fixed recognition laws, and alternating it with updates of the $q_i$ descends one two-argument objective~\cite{Neal1998}; it is the Gaussian restriction of the aggregate-prior proposition of~\cite[Ch.~12]{DennisGaugeVFE2026}. In this model the parent is a fitted hyperparameter: it carries no belief of its own, $q_I$ is a point estimate, not a posterior, and $\Sigma_I$ is the population covariance of the constituents. A different generative model, with one parent latent and the constituents as noisy copies of it (Appendix~\ref{app:augmented_joint}), gives $\Sigma_I$ the meaning of posterior uncertainty; there the mean-field coordinate update of the parent is precision addition, which sharpens the parent as $\mathcal{O}(1/|I|)$, does not depend on the constituent covariances, and runs opposite to the covariance pooling of Eq.~\eqref{eq:meta_agent_sigma_impl}. The two readings are distinct models, and only the random-effects reading makes the barycenter an exact coordinate. The derivation fixes the form and not the weights: any weight vector yields a valid M-step coordinate, so the coherence weight $\exp[-D_{\mathrm{KL}}]$ of Section~\ref{sec:meta_agent_state_construction} is a declared modeling choice.

A further proposal treats the parent state as a predictive sufficient statistic of its constituents in the sense of the information bottleneck~\cite{tishby1999information}: reading $C(I)$ as the bottleneck rate term and the free-energy change as a proxy for the predictive term, the information-bottleneck Lagrangian~\cite{bialek2001predictability} would select the coarse-graining that retains maximal predictive information at fixed rate. This is an analogy, not a derivation. The closed-form Gaussian-IB solution of~\cite{chechik2005information} does not apply directly, because the parent state is a pair of Gaussian-parameterized laws on associated bundles, not a jointly Gaussian vector, and the deterministic~\cite{StrouseSchwab2017} and agglomerative~\cite{SlonimTishby2000} variants are not worked out here.

\subsection{The Threshold Detector: A Coherence Filter for Candidate Clusters}
\label{sec:meta_agent_threshold}

Evaluating~\eqref{eq:meta_agent_FE_criterion} requires the optimized free energy under both configurations, which is expensive, and the criterion is itself conjectural. Candidate clusters are therefore found with a threshold detector built from a coherence proxy that combines post-transport belief alignment, post-transport model alignment, and spatial overlap. Each factor is a bounded exponential of a dispersion, so the detector lives in $[0, 1]$. Belief coherence is the bounded exponential of the average post-transport belief dispersion,
\begin{equation}
C_q(I, x) = \exp\left[-\frac{1}{\tau_q |I|^2}\sum_{i,j \in I} D_{\mathrm{KL}}\big(q_i^{(\ell)}(x) \| \Omega_{ij}[q_j^{(\ell)}](x)\big)\right] \in [0, 1].
\end{equation}
Model coherence is defined analogously as $C_s(I, x) = \exp[-V^{(s)}_I(x)/\tau_s]$ with $V^{(s)}_I(x) = |I|^{-2}\sum_{i,j \in I} D_{\mathrm{KL}}(s_i^{(\ell)}(x) \| \tilde{\Omega}_{ij}[s_j^{(\ell)}](x))$, the superscript labeling the model channel and the subscript the cluster. Presence measures spatial overlap via $P(I, x) = |I|^{-1}\sum_{i \in I} \chi_i(x) \in [0, 1]$. The dispersion temperatures $\tau_q, \tau_s > 0$ set the belief and model resolution at which coherence is read off. The consensus score is the product $\Gamma(I, x) = P(I, x) C_q(I, x) C_s(I, x) \in [0, 1]$, and a cluster is a candidate when $\Gamma(I, x) > \Gamma_{\min}$ and $|I| \geq N_{\min}$. Boundedness matters because $D_{\mathrm{KL}}$ is unbounded above, so a $1 - D_{\mathrm{KL}}$ surrogate would be signed and could yield misleadingly positive products from two negative factors. The same exponential detector reappears in Appendix~\ref{app:rigorous_rg}, written with the barycenter-relative dispersions $V_I^{(q)}$, $V_I^{(s)}$ and strengthened by a spectral-gap condition. Numerical values quoted below use $\Gamma_{\min} = 1/2$, $N_{\min} = 2$, and $\tau_q = 0.05$ for illustration; they are not theoretical requirements.

The detector's relation to the licensing criterion is partial. Because $P, C_s \le 1$, the pass band $\Gamma > \Gamma_{\min}$ forces the mean belief dispersion below $\tau_q\log(1/\Gamma_{\min})$, which selects the high-coherence end, including the limit of vanishing dispersion. The criterion has no derived optimum against which this band could be compared (Section~\ref{sec:meta_agent_variational}), and we do not establish that the product $P \cdot C_q \cdot C_s$ tracks any free-energy change, even in the high-coherence limit. The strengthened detector of Appendix~\ref{app:rigorous_rg} additionally requires a positive constrained spectral gap $m_I > m_{\min}$; the canonical specification above does not check this condition, so by the appendix's own classification it operates as a similarity threshold and not as a closure surrogate. The appendix's retention theorem (Theorem~\ref{thm:rg_detector_retention}) concerns the strengthened detector with the gap condition, not the one specified here.

\subsection{Gauge-Covariant Belief Construction}
\label{sec:meta_agent_state_construction}

When a meta-agent $I$ forms, its belief is the barycenter of~\eqref{eq:meta_agent_barycenter}. The working form drops the dispersion term in~\eqref{eq:meta_agent_sigma_barycenter}, an approximation for high-coherence clusters. The dropped term is the constituent mean-spread $\sum_i w_i (\tilde{\mu}_i - \bar{\mu}_I)(\tilde{\mu}_i - \bar{\mu}_I)^\top$. With normalized weights and $\bar{\varepsilon}_w = \sum_{i,j} w_i w_j D_{\mathrm{KL}}(q_i \| \Omega_{ij}[q_j])$, it obeys the operator-norm bound
\begin{equation}
\Big\|\sum_i w_i (\tilde{\mu}_i - \bar{\mu}_I)(\tilde{\mu}_i - \bar{\mu}_I)^\top\Big\|_{\mathrm{op}} \le \max_j \lambda_{\max}(\tilde{\Sigma}_j) \bar{\varepsilon}_w,
\label{eq:dispersion_drop_bound}
\end{equation}
which follows from the pairwise-difference identity for weighted spreads, $\sum_i w_i (\tilde{\mu}_i - \bar{\mu}_I)(\tilde{\mu}_i - \bar{\mu}_I)^\top = \tfrac12\sum_{i,j} w_i w_j (\tilde{\mu}_i - \tilde{\mu}_j)(\tilde{\mu}_i - \tilde{\mu}_j)^\top$, the Gaussian mean-term lower bound $D_{\mathrm{KL}}(\mathcal{N}_1 \| \mathcal{N}_2) \ge \tfrac12(\Delta\mu)^\top\Sigma_2^{-1}(\Delta\mu)$, and cocycle invariance of the transported divergences, which holds in Regime I (trivial holonomy). The bound is tight for collinear equal-covariance configurations. The detector controls the unweighted mean dispersion $\bar{\varepsilon} = |I|^{-2}\sum_{i,j} D_{\mathrm{KL}}(q_i \| \Omega_{ij}[q_j])$, not $\bar{\varepsilon}_w$, and the two are related by $\bar{\varepsilon}_w \le |I|^2 (\max_i w_i)^2 \bar{\varepsilon}$. Inside the pass band the right side of~\eqref{eq:dispersion_drop_bound} is therefore at most $|I|^2(\max_i w_i)^2 \tau_q\log(1/\Gamma_{\min}) \max_j \lambda_{\max}(\tilde{\Sigma}_j)$, which reduces to $\tau_q \log 2 \max_j \lambda_{\max}(\tilde{\Sigma}_j) \approx 0.035 \max_j \lambda_{\max}(\tilde{\Sigma}_j)$ for uniform weights at the illustrative values. For uniform weights the truncation is thus a moment-matching approximation with a controlled error; for nonuniform weights the factor $|I|^2(\max_i w_i)^2$ must be carried, and at order-one dispersion temperatures the mean-dispersion bound degrades to $\log 2 \approx 0.69$ nats, a regime in which the dispersion term should be reinstated. The working forms keep the barycenter mean and the transported-covariance pooling and drop only the constituent mean-spread:
\begin{align}
\bar{\mu}_I(x) &= \frac{\sum_{i \in I} w_i(x) \Omega_{Ii}[\mu_i^{(\ell)}](x)}{\sum_{i \in I} w_i(x)},
\label{eq:meta_agent_mu_impl}
\\
\bar{\Sigma}_I(x) &= \frac{\sum_{i \in I} w_i(x) \Omega_{Ii}[\Sigma_i^{(\ell)}](x) \Omega_{Ii}^\top}{\sum_{i \in I} w_i(x)}.
\label{eq:meta_agent_sigma_impl}
\end{align}
The weights depend on presence and coherence through $w_i^I(x) = \chi_i(x) \exp[-D_{\mathrm{KL}}(\Omega_{Ii}[q_i^{(\ell)}] \| \bar{q}_I^{(\ell)})]$, with the comparison taken after transport into the meta-agent's frame. The scalar weight is gauge invariant, being built from a transported divergence and a support mask that are each invariant under common pushforward, as the covariance theorem of Appendix~\ref{app:rigorous_rg} requires. It is a declared modeling choice (Section~\ref{sec:meta_agent_variational}), not the approximation of an optimal weight. Because $\bar{q}_I^{(\ell)}$ appears on the right-hand side through the weights, the formula is implicit and is solved by fixed-point iteration. Its convergence is guaranteed when the maximum pairwise transported divergence satisfies $\varepsilon < \varepsilon_0$ (Lemma~\ref{lem:weight_contraction}). The detector's pass band bounds only the mean dispersion, which implies $\max_{i \ne j} D_{\mathrm{KL}}(q_i \| \Omega_{ij}[q_j]) \le |I|^2\tau_q\log(1/\Gamma_{\min})$ and not $\varepsilon < \varepsilon_0$, so for large clusters a separate check of the maximum divergence is required. The check is not a formality: with identity covariances, 609 coincident members and one member at distance $4.6$ have mean dispersion $0.0346 < \tau_q\log 2$ at $\tau_q = 0.05$, maximum pairwise divergence $10.6$, and a weight map with two stable fixed points, one of which gives the outlier $0.974$ of the total weight. A non-convergence test cannot detect this bistability, because the iteration converges.

\begin{lemma}[Contraction of the weight map at high coherence]
\label{lem:weight_contraction}
Fix $0 < \lambda_{\min} \le \lambda_{\max}$, assume Regime I transport, and suppose the transported constituents satisfy $\lambda_{\min} I \preceq \tilde{\Sigma}_i \preceq \lambda_{\max} I$, the supports satisfy $\sum_{i \in I}\chi_i > 0$, and the pairwise transported divergences satisfy $\max_{i \ne j} D_{\mathrm{KL}}(q_i \| \Omega_{ij}[q_j]) \le \varepsilon$. Let $T$ send a Gaussian $\bar{q} = \mathcal{N}(\bar{\mu}, \bar{\Sigma})$ with $\lambda_{\min} I \preceq \bar{\Sigma} \preceq \lambda_{\max} I$ to the pair of Eqs.~\eqref{eq:meta_agent_mu_impl}--\eqref{eq:meta_agent_sigma_impl} with weights $w_i(\bar{q}) = \chi_i \exp[-D_{\mathrm{KL}}(\Omega_{Ii}[q_i] \| \bar{q})]$. In the coordinates $(\mu, \Sigma)$ with the Euclidean and Frobenius norms, $T$ maps this bounded-spectrum set into the convex hull $H$ of the moments $\{(\tilde{\mu}_i, \tilde{\Sigma}_i)\}$, and on $H$ it is Lipschitz with constant $C(\lambda_{\min}, \lambda_{\max})\varepsilon$. For $\varepsilon < \varepsilon_0 = 1/C$ the fixed point is unique and the iteration converges geometrically from any start in the bounded-spectrum set. Every admissible constant satisfies $\varepsilon_0 \le 2$ when $\lambda_{\min} = \lambda_{\max} = 1$.
\end{lemma}

\begin{proof}[Proof sketch]
With normalized weights $\bar{w}_i$, $T$ is the convex combination $\sum_i \bar{w}_i \zeta_i$ of the moments $\zeta_i = (\tilde{\mu}_i, \tilde{\Sigma}_i)$ (not of the expectation parameters, which would reinstate the dispersion term), so it maps into $H$, and $H$ lies in the bounded-spectrum set. Writing $g_i$ for the gradient of $D_{\mathrm{KL}}(\Omega_{Ii}[q_i] \| \bar{q})$ in $(\bar{\mu}, \bar{\Sigma})$ and $\bar{g} = \sum_i \bar{w}_i g_i$, the derivative is $dT[h] = -\sum_i \bar{w}_i \langle g_i - \bar{g}, h\rangle (\zeta_i - \bar{\zeta})$, a weighted covariance between moment differences and gradient differences. By cocycle invariance, $D_{\mathrm{KL}}(\Omega_{Ii}[q_i] \| \Omega_{Ij}[q_j]) = D_{\mathrm{KL}}(q_i \| \Omega_{ij}[q_j]) \le \varepsilon$, so the Gaussian divergence bounds on the bounded-spectrum set give $\|\zeta_i - \zeta_j\| = \mathcal{O}(\sqrt{\varepsilon})$. The mean part of $g_i - g_j$ is $\bar{\Sigma}^{-1}(\tilde{\mu}_j - \tilde{\mu}_i)$, independent of $\bar{q}$, and the covariance part is $\mathcal{O}(\sqrt{\varepsilon})$ on the bounded set $H$; hence $\|dT\| = \mathcal{O}(\varepsilon)$ on the convex set $H$, and the Banach fixed-point theorem applies there. For the bound on $\varepsilon_0$, take two members $\mathcal{N}(\pm a, 1)$ in one dimension, with pairwise divergence $\varepsilon = 2a^2$. After one step $\bar{\Sigma} = 1$, and the mean map is $m \mapsto a\tanh(am)$, with slope $\varepsilon/2$ at $m = 0$; for $\varepsilon > 2$ it has three fixed points, so uniqueness fails.
\end{proof}

The defining property of the construction is gauge consistency: all constituent beliefs are first transported to the meta-agent's frame via $\Omega_{Ii}$ before averaging, so the scalar weights are gauge invariant while the averaged belief $\bar{q}_I$ transforms by pushforward with the meta-agent's frame. Without transport into a common frame the average would depend on the agents' passive choices of frame. By Proposition~\ref{prop:parent_frame_gauge} the parent frame cancels from every constituent law, so the licensing argument inherits no frame error when one convention for $U_I$ is used consistently.

The meta-agent's gauge frame and its independent model-frame analog are, in the working form, the Lie-algebra-additive averages
\begin{equation}
\begin{aligned}
\phi_I^{(\ell+1)}(x)
&= \frac{\sum_{i \in I} w_i(x)\phi_i^{(\ell)}(x)}{\sum_{i \in I} w_i(x)}, \\
\tilde\phi_I^{(\ell+1)}(x)
&= \frac{\sum_{i \in I} w_i^{\mathrm{mod}}(x)\tilde\phi_i^{(\ell)}(x)}{\sum_{i \in I} w_i^{\mathrm{mod}}(x)}.
\end{aligned}
\end{equation}
The first is the first fixed-point iterate, from the identity, of the exponential barycenter $\sum_i w_i\log(U_I^{-1}U_i) = 0$~\cite{PennecArsigny2013}: the iteration $U \mapsto U\exp[\sum_i w_i \log(U^{-1}U_i)]$ started at $U = I$ returns $\exp(\sum_i w_i \phi_i)$ when $U_i = \exp\phi_i$. It coincides with the exponential barycenter for commuting $\phi_i$, and in particular for abelian $G$~\cite{Hall2015}. It is not gauge-equivariant: its defect under $U_i \mapsto gU_i$ is controlled by the magnitude of the gauge transformation and by the spread of the $\phi_i$, not by $\|\phi_i\|$. The model-fiber weight $w^{\mathrm{mod}}_i$ is built from model-fiber transported divergences and solved by an independent fixed point on the model fiber, so the cross-scale model transport $\tilde\Omega_{iI}$ is well-defined at meta levels; in the tied-frame configuration $\tilde\phi_I$ is the $\rho_{\mathrm{model}}$-image of $\phi_I$ and no separate average is required. Reinstating the dispersion term in $\bar{\Sigma}_I$ recovers the strict moment-matched mixture covariance at modest computational cost.

\subsection{The Consensus Foil: Pointwise Consensus}
\label{sec:epistemic_collapse}

One limiting case brackets the construction and clarifies what it is not. In pointwise consensus, transported beliefs and models coincide, $\Omega_{ij}[q_j](c) = q_i(c)$ and $\tilde{\Omega}_{ij}[s_j](c) = s_i(c)$ for all $c$ in the shared support. Every inter-agent coupling term then attains its minimum value zero on both channels, and the couplings exert no force on beliefs or models, so learning is driven only by observations. Pointwise consensus is not the criterion for meta-agent formation, and applied across all agents it would be a static state in which the agents' distinct perspectives no longer register in the coupling terms. Formation requires bounded, not vanishing, internal disagreement modulo transport, with closure relative to the outside. Pointwise consensus nonetheless has a constructive role at the boundary: constituents that have collapsed to common values up to transport can be replaced by a single representative without information loss. The construction requires only that the system as a whole not reach pointwise consensus, from which it is kept away by continual observation and by hierarchical feedback that perturbs constituents heterogeneously.

